Incomplete telemetry complicates microservice resource forecasting and can change the accuracy gained from more complex models. We examine how temporal decomposition, availability descriptors and model size affect CPU and memory forecasts across five experimental phases. The historical ChronoGraph experiment uses 28 services, 8005 timestamps, a chronological 70/10/20 split, twelve-step histories, three training seeds and twenty paired masks. Experiments with E0--E5 and compact C0--C3/V2 motivate calibrated transfer to 32 additional services. The expanded comparison evaluates fifteen families with ten primary neural seeds and twenty masks. Secondary experiments vary missingness, direct horizons and history lengths. Frozen models also transfer to one day of Alibaba 2022 after six hours of calibration. Historically, E5 reduces paired memory MASE by 1.280\% relative to DLinear and increases CPU MASE by 0.984\%. V2 does not retain the full memory gain. On the additional services, A3 uses 15 parameters and improves memory MASE by 0.02523 relative to E5's 74-parameter configuration. All ten seed means favor A3 (Holm-adjusted exact sign-test $p=0.0078125$), with improvement in only 10 of 32 services. A3 does not generally improve on DLinear. Small PatchTST and NLinear have the lowest observed primary MASE means, while persistence ranks first externally. A3 achieves CPU/CUDA median forward times of 0.050/0.141 ms with resident inputs for one origin and 32 services. The results establish context-specific tradeoffs under synthetic loss, shared ChronoGraph periods and short external coverage.
